
\documentclass[11pt,a4paper]{report}
\usepackage{ai_security_bluebook}

% =============================================================================
% Chapter--pillar mapping (for navigation; not printed)
% Ch1: Unified AI Security Control Plane (primary pillar: \ControlPlane{})
% Ch2: Decoupled Safety (primary pillar: \DecoupledSafety{})
% Ch3: Reference architecture + \SafetyKernel{} (Control Plane + Decoupled Safety)
% Ch4: Runtime enforcement (Control Plane + Decoupled Safety)
% Ch5: Evaluation framework (primary pillar: \EvalFramework{})
% Ch6: Governance + roadmap (all three pillars; evidence and org closure)
% =============================================================================

% =============================
% Metadata
% =============================
\bluebooktitle{Blueprint for Enterprise AI Security}
\bluebooksubtitle{Unified AI Security Control Plane, Decoupled Safety, and Evaluation Framework}
\bluebookauthor{Yunhao Wang et al.}
\bluebookorg{Digital Trust / Tianmu / AI Security Research}
\bluebookversion{v1.0-draft}
\bluebookdate{\today}
\bluebookclassification{Internal Draft}
\bluebookkeywords{AI Security; Unified Control Plane; Decoupled Safety; Runtime Enforcement; Evaluation Framework}
\bluebookcovernote{This technical bluebook specifies a unified control-plane architecture, a decoupled safety model, and an evidence-centric evaluation methodology for enterprise AI systems. Recommended compiler: XeLaTeX.}

\begin{document}
\makebluebookfrontmatter

\begin{bluebookabstract}
This bluebook defines a unified architectural and evaluation blueprint for enterprise-grade AI security. It organizes AI security around three mutually reinforcing pillars: (i) a \ControlPlane{} that externalizes and orchestrates security decisions across the AI lifecycle; (ii) a \DecoupledSafety{} paradigm that treats safety as an explicit, compositional control layer rather than a hidden property of any single model; and (iii) a rigorous \EvalFramework{} that measures not only attack resistance, but also controllability, utility retention, auditability, and operational sustainability. The central claim is that enterprise AI security is a \emph{control-plane} problem: assurance must be produced by composable runtime enforcement tied to measurable evidence, not assumed from static model alignment alone.
\end{bluebookabstract}

\executivesummary{
\begin{itemize}
  \item \textbf{Context.} Production AI is shifting from isolated model endpoints to compound systems: orchestration, retrieval-augmented generation (\RAG{}), tools, memories, and multi-turn agents. Failure modes span confidentiality, integrity, availability, and operational governance---not only toxic outputs.
  \item \textbf{Problem.} Capability growth has outpaced enterprise-grade mechanisms for consistent policy enforcement, audit traceability, and release-grade assurance under non-determinism and continuous change.
  \item \textbf{Thesis.} Security and safety constraints should be implemented as an externalized, composable, auditable \ControlPlane{} with a \SafetyKernel{} that enforces policy across inputs, retrieval, inference, outputs, and tool calls. Model alignment is necessary but not sufficient; \DecoupledSafety{} makes obligations explicit and testable.
  \item \textbf{Evidence thesis.} Claims require an \EvalFramework{} that binds benchmarks, acceptance contracts, online telemetry, red teaming, regression, and release artifacts into a single evidentiary chain.
  \item \textbf{Deliverables.} A reference threat model and objectives; a layered reference architecture; runtime policy orchestration patterns including \RAG{} and tool-use controls; metric matrices for security, utility, system cost, and governance; a phased deployment and research roadmap; and appendices for jurisdiction-aware evaluation patterns and terminology.
\end{itemize}
}

% =============================
% Chapter 1
% =============================
\chapterstarter{Unified AI Security Control Plane: Problem Statement and System Boundary}{
This chapter states the enterprise AI security problem as a control-plane problem. \textbf{Pillar focus:} \ControlPlane{}. It establishes strategic context, a formal boundary between judgment and execution, adversary assumptions, and measurable security objectives including confidentiality, integrity, availability alongside controllability and auditability.
}{
\begin{itemize}
  \item Problem statement and scope; system boundary and trust assumptions.
  \item Structured threat model aligned to runtime control points.
  \item Security objectives (CIA plus controllability, auditability, composability, utility preservation) and explicit trade-offs.
\end{itemize}
}

\section{Strategic context and enterprise shift}
\textbf{Pillar:} \ControlPlane{}. Enterprise AI deployments increasingly behave as distributed systems: gateways, routers, retrieval indices, vector stores, workflow engines, plugins, and human workflows. Risk is therefore \emph{systemic}: prompt injection crosses trust boundaries; poisoned corpora manipulate answers; tools chain into privileged actions; orchestration state creates cross-turn attacks; logging gaps erase accountability. Traditional application and content-safety tooling, when used as disconnected add-ons, does not provide a single enforcement and evidence standard across these paths.

\section{System boundary and responsibility split}
\textbf{Pillar:} \ControlPlane{}. This bluebook distinguishes (i) \emph{security judgment}: policy interpretation, risk scoring, allow/deny/rewrite decisions, explainability payloads, and evidence emission; from (ii) \emph{execution platform}: networking, scaling, orchestration \textsc{sla}, and workload placement. The \ControlPlane{} must remain semantically anchored in policy and evidence rather than collapsing into either a model-only narrative or an undifferentiated infrastructure story. External context (documents, tools, APIs) is part of the attack surface and must fall under the same envelope of controls and traceability.

\begin{problemstatement}
The central question is not merely how to make a model safe in isolation, but how to guarantee that an \LLM{}-centric system remains within explicit, auditable risk boundaries under dynamic context injection, external retrieval, tool invocation, memory and state, and heterogeneous deployment constraints. The solution is not ``more detectors'' disconnected from measurement and release, but a coherent control plane: unified semantics for risk, deterministic enforcement hooks, and an evaluation system that produces release-grade evidence.
\end{problemstatement}

\section{Assets, adversary model, and assumptions}
\textbf{Pillar:} \ControlPlane{} / \EvalFramework{}. The assets include user and enterprise data confidentiality, integrity of prompts and retrieved content, integrity and non-repudiation of tool side effects, availability of the service within \textsc{sla} envelopes, and integrity of audit records. We consider adversaries who can craft adversarial prompts, poison or leak via retrieval channels, misuse or chain tool calls, manipulate session and workflow state, and abuse operator surfaces (e.g., shadow prompts in documents). Insider threat is addressed primarily through separation of duties, evidence immutability, and least privilege on tooling---not through model alignment alone.

Operational assumptions include: policies can be stated in machine-checkable fragments; enforcement points exist at gateway, retrieval, inference, tool, and logging layers; some residual risk is accepted explicitly through acceptance contracts; and non-determinism must be bounded by statistical test plans and online monitoring rather than dismissed.

\begin{threatmodel}
\begin{threattable}
Prompt interface & Policy bypass / exfiltration via instructions & Jailbreaks, delimiter attacks, indirect prompt injection & Policy failure, data disclosure, privilege abuse & Input governance, normalization, instruction boundaries \\
Retrieval path & Corpus manipulation / sensitive exposure & Poisoning, conflicting sources, injected instructions in text & Wrong or harmful answers, leakage, compliance breach & Retrieval security, provenance, segmentation, cite binding \\
Tool-use layer & Escalation of privilege & Tool chain attacks, forged parameters, ambiguous approvals & Data loss, financial or ops impact, lateral movement & Tool policies, approvals, sandboxing, rate limits \\
Model execution & Unsafe generation under benign intent & Misalignment bursts, hallucinated facts & Harmful content, unreliable automation & Output governance, grounding constraints, calibrated abstention \\
Runtime orchestration & Control-plane inconsistency & Cross-turn state manipulation, stale policy binding & Silent policy drift, audit gaps & Policy engine versioning, atomic decision logs \\
\end{threattable}
\end{threatmodel}

\section{Security objectives and trade-offs}
\textbf{Pillar:} \ControlPlane{} / \EvalFramework{}. Objectives must be simultaneously satisfied at the system level; trade-offs are explicit and measurable.

\begin{securityobjectives}
\begin{itemize}
  \item \textbf{Confidentiality.} Sensitive enterprise and user data must not cross trust boundaries via prompts, retrieval, logs, or tool outputs beyond policy.
  \item \textbf{Integrity.} Prompts, retrieved passages, tool parameters, and side effects must be integrity-protected within threat assumptions; tampering must be detectable in audit evidence.
  \item \textbf{Availability.} Security controls must preserve service viability within latency and throughput envelopes; degradation must be graceful and measurable.
  \item \textbf{Controllability.} High-risk behaviors must be constrainable through explicit runtime policy rather than implicit model habit.
  \item \textbf{Auditability.} Policy-relevant decisions emit structured evidence suitable for replay, governance review, and regression testing.
  \item \textbf{Composability.} Controls remain meaningful across heterogeneous models, revisions, and stacks without tight coupling to any single checkpoint file.
  \item \textbf{Utility preservation.} Security interventions minimize unnecessary harm to legitimate task success; unnecessary blocks are treated as defects with metrics.
\end{itemize}
\end{securityobjectives}

\noindent\textbf{Trade-offs.} Strong blocking improves integrity and confidentiality metrics but may reduce utility and availability; retrieval segmentation improves confidentiality but may reduce answer completeness; audit depth improves governance metrics but consumes storage and latency budget. \DecoupledSafety{} exists precisely to negotiate these trade-offs \emph{as policy}, not as ad hoc tuning.

\designprinciple{Decoupling}{Security logic is externalized from any single model and expressed through explicit policy interfaces enforced by a \SafetyKernel{}.}
\designprinciple{Cross-layer enforcement}{Controls span input, retrieval, inference, tool-use, output, and governance logging as one coherent envelope.}
\designprinciple{Evidence before trust}{Every assurance claim is backed by tests, telemetry, or audit artifacts suitable for release decisions.}

\keytakeaway{The \ControlPlane{} is the architectural unit that makes enterprise AI security operationally enforceable, model-agnostic, and evolvable under changing stacks and policies.}

% =============================
% Chapter 2
% =============================
\chapterstarter{Decoupled Safety: Conceptual Foundation and Security Model}{
This chapter formalizes \DecoupledSafety{}: safety as an externalized, compositional layer with explicit precedence over model preferences. \textbf{Pillar focus:} \DecoupledSafety{}. It defines concepts, invariants, and characteristic failure modes when safety is mistaken for intrinsic model properties.
}{
\begin{itemize}
  \item Definitions and invariants for policy precedence and composability.
  \item Failure modes when alignment is treated as sufficient for system safety.
  \item Bridge from conceptual model to kernel-level enforcement (preview for Chapter~3).
\end{itemize}
}

\section{Limits of model-only safety narratives}
\textbf{Pillar:} \DecoupledSafety{}. Alignment and instruction tuning reduce baseline harmful completion rates but do not establish system-level guarantees under retrieval, tools, adversarial contexts, or orchestration. A model may appear ``safe'' in offline evaluations yet fail when connected to corpora that carry adversarial instructions, tools with real side effects, or workflows that accumulate unsafe state across turns.

\begin{definition}[Decoupled Safety]
\DecoupledSafety{} is the design principle that safety constraints are represented, enforced, and audited in a dedicated control layer that remains separable from the internal parameters of any individual model. The policy layer exposes stable interfaces (events, obligations, decisions) independent of provider-specific checkpoints.
\end{definition}

\begin{invariant}[Policy precedence]
Whenever a conflict occurs between generative preference and explicit safety policy, the policy layer must dominate the final system behavior. Partial outcomes (e.g., rewritten or redacted outputs) must still satisfy policy.
\end{invariant}

\begin{invariant}[Compositional black-box interfaces]
Safety enforcement consumes typed \emph{security events} (inputs, retrieved units, tool invocations, candidate outputs) and emits \emph{decisions} and \emph{evidence records}. Downstream modules must not bypass these decisions without an explicit, auditable escalation path.
\end{invariant}

\section{Security model abstractions}
\textbf{Pillar:} \DecoupledSafety{}. For implementation and evaluation, separate \emph{policy}~$P$ (obligations and thresholds), \emph{enforcement kernel}~$K$ (guards, routers, tool gate, orchestration hooks), \emph{environment}~$E$ (deployment topology and trust zones), and \emph{evidence}~$\mathrm{Ev}$ (auditable traces). Changing the model checkpoint should chiefly alter generation statistics; it must not silently invalidate $P$ without corresponding evidence updates.

\section{Characteristic failure modes}
\securityrisk{Ambiguous ownership}{When safety logic hides inside prompts or model weights, upgrades silently rewrite effective policy without regression detection.}
\securityrisk{Orchestration bypass}{Agents and workflows route around guards using alternate tools or implicit channels; consistency fails without orchestration-bound enforcement.}
\securityrisk{Evidence absence}{Operators cannot reconstruct \emph{why} a decision occurred, blocking audits, postmortems, and supervised learning from incidents.}

\researchagenda{Foundational question}{Under which conditions can compositional safety constraints be enforced with measurable error bounds across heterogeneous models, non-deterministic decoding, and multi-step tool use?}

\keytakeaway{\DecoupledSafety{} converts safety from an emergent property of training into an explicit engineering object: policy, kernel, evidence, and evaluation.}

% =============================
% Chapter 3
% =============================
\chapterstarter{Reference Architecture and Decoupled Safety Kernel}{
This chapter specifies a reference architecture for the \ControlPlane{} and decomposes the \SafetyKernel{}. \textbf{Pillar focus:} \ControlPlane{} + \DecoupledSafety{}. It describes layered responsibilities and the parallel flows of data, control, and evidence required for auditability.
}{
\begin{itemize}
  \item Layered architecture from interaction through governance.
  \item Module interfaces for policy engine, guards, retrieval mediation, tool control, audit plane.
  \item Initial deliverables and phased ownership (see also Chapter~6).
\end{itemize}
}

\section{Layered reference architecture}
\textbf{Pillar:} \ControlPlane{}. The architecture is organized into layers with crisp responsibilities:
\begin{itemize}
  \item \textbf{Interaction layer.} User and application inputs, authentication and tenancy context, session management, request shaping.
  \item \textbf{Context layer.} Prompt understanding, routing, retrieval planning, provenance binding, segmentation of corpora by sensitivity and purpose.
  \item \textbf{Model execution layer.} Model routing, decoding controls, grounding hooks, abstention and uncertainty signaling where available.
  \item \textbf{Runtime enforcement layer.} Policy engine, guardrails, retrieval mediation, response filtering, tool-use gate, deterministic rewrite and redaction, escalation.
  \item \textbf{Governance layer.} Evidence stores, evaluation pipelines, release gates, monitoring, incident replay, policy lifecycle management.
\end{itemize}

\section{Data plane, control plane, and evidence plane}
\textbf{Pillar:} \ControlPlane{}. The \emph{data plane} carries prompts, retrieved texts, embeddings metadata, tool requests and responses, and completions. The \emph{control plane} carries policies, decisions, obligations, routing outcomes, risk scores, and escalation states. The \emph{evidence plane} binds identifiers across systems: policy version, model route, retrieval corpus slice, detector versions, hashes of prompts and documents where appropriate, decisions, and operator approvals. Assurance arguments depend on correlating evidence records across planes.

\section{\SafetyKernel{} interfaces}
\textbf{Pillar:} \DecoupledSafety{}. The \SafetyKernel{} exposes narrow, stable interfaces so enforcement remains composable:
\begin{itemize}
  \item \textbf{Event ingestion.} Normalized events for inputs, retrieval hits, candidate completions, tool calls.
  \item \textbf{Policy evaluation.} Deterministic evaluation on defined events with versioned rule sets and thresholds.
  \item \textbf{Action dispatch.} Allow, rewrite, constrain, sandbox, escalate, block, degrade with typed parameters.
  \item \textbf{Evidence emission.} Structured records consumable by \EvalFramework{} and governance tooling.
\end{itemize}

\begin{deliverabletable}
Policy engine & Schema + rules + decision traces & Deterministic replay for defined event types & Runtime policy & P1 \\
Input / output \Guardrails{} & Classifiers, validators, rewrite templates & Reduced measured bypass vs.\ baselines & Security ML & P1 \\
Retrieval security service & Corpus labeling, cite binding, retrieval allow lists & Reduced successful injection/poisoning impact & Platform + App & P1 \\
Tool controller & Capability matrix, consent, approvals, sandbox & High-risk actions require explicit policy paths & Agent platform & P2 \\
Audit plane & Append-only logs, correlation identifiers & Full incident reconstruction from evidence & Platform SRE & P2 \\
Evaluation services & Suites, dashboards, regressions & Release gates satisfied with versioned bundles & Assurance & P1 \\
\end{deliverabletable}

\keytakeaway{A \SafetyKernel{} turns \DecoupledSafety{} into implementable interfaces: typed events, deterministic decisions, and evidence suitable for governance.}

% =============================
% Chapter 4
% =============================
\chapterstarter{Runtime Enforcement and Policy Orchestration}{
This chapter specifies how policies are enforced end-to-end, including risk-adaptive control selection, retrieval and tool mediation, and human-in-the-loop escalation. \textbf{Pillar focus:} \ControlPlane{} + \DecoupledSafety{}. It connects the conceptual model to measurable runtime behavior.
}{
\begin{itemize}
  \item Risk signals, policy binding, interventions, logging, and feedback to evaluation.
  \item Concrete patterns for \RAG{} and tool-use enforcement and safe degradation.
\end{itemize}
}

\section{Risk-adaptive control loop}
\textbf{Pillar:} \ControlPlane{}. Online enforcement follows a closed loop aligned with assurance:
\begin{enumerate}
  \item \textbf{Measure.} Offline and online measurements produce decision signals (attack success rates, drift, refusal quality, latency).
  \item \textbf{Profile.} Per-model and per-scenario risk profiles summarize exposure given deployment context.
  \item \textbf{Select.} Choose control depth (rules vs.\ lightweight detectors vs.\ heavier checks; tool restrictions; retrieval segmentation) subject to risk, cost, latency, and business tolerance.
  \item \textbf{Enforce.} Apply actions deterministically through the \SafetyKernel{} interfaces.
  \item \textbf{Evidence and feedback.} Emit audit records; failures feed benchmark updates, red-team priorities, and policy revisions.
\end{enumerate}

\section{Policy decision lifecycle}
\textbf{Pillar:} \DecoupledSafety{}. At request time:
\begin{enumerate}
  \item Inspect the request and derive risk signals (topic sensitivity, channel trust, retrieval scope, tool exposure).
  \item Bind the request to applicable policy packs and versions (tenant, geography, workload class).
  \item Evaluate interventions: allow, rewrite, constrain, sandbox, escalate, block, degrade with explicit rationale codes.
  \item Log decisions with correlated evidence payloads suitable for replay.
  \item Feed incidents and near-misses into offline suites and shadow evaluations.
\end{enumerate}

\section{Retrieval security (\RAG{})}
\textbf{Pillar:} \ControlPlane{}. Retrieval introduces third-party text into the prompt boundary; treat retrieved segments as untrusted code-like content. Controls include corpus segmentation (public vs.\ confidential), metadata and lineage requirements, retrieval-time filtering, citation binding to passages, contradiction handling, maximum privilege scopes for answers derived from sensitive slices, and poison detection hooks (batch and online). Retrieval policies must co-vary with evaluation scenarios that include poisoned-document suites.

\section{Tool-use control}
\textbf{Pillar:} \ControlPlane{}. Tools map language to side effects; enforce capability matrices per tenant and workflow, typed arguments validation, dynamic approval for high-impact actions, sandboxed execution where feasible, durable rate limits, and separation between read-only vs.\ mutating tools. Orchestrators must not construct parallel paths that bypass recorded decisions.

\section{Escalation, fallback, and human-in-the-loop}
\textbf{Pillar:} \DecoupledSafety{}. When automated confidence is insufficient, escalation routes include constrained execution (safe templates), partial answers, human approval queues with \textsc{sla}, and incident capture for training and policy updates. Fallback strategies are themselves policy-governed to avoid silent degradation of security posture.

\evaluationnote{Runtime assurance target}{Each intervention path must have paired metrics: residual risk, utility impact, latency delta, and audit completeness.}

\keytakeaway{Runtime enforcement realizes \DecoupledSafety{} as measurable decisions on real events, not as static configuration files divorced from evidence.}

% =============================
% Chapter 5
% =============================
\chapterstarter{Evaluation Framework: Security, Utility, and System Assurance}{
This chapter defines the \EvalFramework{} as the evidence engine for the \ControlPlane{}. \textbf{Pillar focus:} \EvalFramework{}. It specifies semantics for benchmarks, acceptance contracts, metrics across security and utility, multi-track evaluation under shared infrastructure, and release-grade evidence packs.
}{
\begin{itemize}
  \item Metric categories and mapping from capabilities to tests.
  \item Offline, shadow, adversarial, and regression workflows tied to release gates.
  \item Governance metrics for audit completeness and policy conformance.
\end{itemize}
}

\section{Semantics, benchmarks, and acceptance contracts}
\textbf{Pillar:} \EvalFramework{}. Risks must be classified under a unified taxonomy linking detection categories, scoring rubrics, and obligations. A \emph{benchmark constitution} states coverage expectations (what classes of failures must be measured and how aggregates are computed). An \emph{acceptance contract} binds scenarios to thresholds and exceptions, translating measurements into admissibility for deployment stages. Contracts are versioned jointly with policies and detector bundles.

\section{Metric layers}
\textbf{Pillar:} \EvalFramework{}. Evaluate across four complementary layers:
\begin{itemize}
  \item \textbf{Security.} Attack success probability, bypass rates, injection impact, unsafe tool invocation, exfiltration attempts.
  \item \textbf{Utility.} Task success, correctness where labeled, helpfulness proxies, calibrated refusal behavior, over-refusal rates where applicable.
  \item \textbf{System.} Latency overhead, throughput impact, memory and compute costs, stability under load.
  \item \textbf{Governance.} Trace completeness, replay fidelity, immutability checks, policy binding correctness, dashboard drill-down for incidents.
\end{itemize}

\section{Multi-track evaluation under shared infrastructure}
\textbf{Pillar:} \EvalFramework{}. Distinct regulatory or organizational obligations need not collapse into a single score. A practical pattern is a \emph{shared execution backbone} (run definitions, targets, artifacts, result stores, \textsc{ci}/\textsc{cd} hooks) with \emph{multiple evaluation tracks} that differ in suites, aggregation logic, and report templates---while converging on a unified decision board for release. Jurisdiction-specific tracks (for example, regional regulatory categories with localized case libraries) should be first-class rather than thin translations of a global suite; evidence packs reference track identifiers explicitly.

\begin{evaluationcriteria}
\begin{evaluationmatrixtable}
Input / IO \Guardrails{} & Bypass rate, false negative on harms & Task completion, user friction & Added latency, tail latency & Red-team prompts, adaptive attacks \\
Retrieval control & Poison hit rate, unauthorized slice access & Answer correctness / usefulness & Retrieval latency, index cost & Poisoned corpora, injection docs \\
Tool-use control & Unsafe action rate, privilege misuse & Workflow success & Approval latency & Agent attack scenarios, chained tools \\
Dual-track / regional suites & Track-specific pass thresholds & Comparable utility baselines & Run cost, operator time & Localized suites + shared harness \\
Audit / evidence plane & Trace completeness, replay success & Debug time to root cause & Storage, log volume & Incident replay drills, sampling audits \\
\end{evaluationmatrixtable}
\end{evaluationcriteria}

\section{Protocols: offline, online, regression, and release}
\textbf{Pillar:} \EvalFramework{}. \textbf{Offline} suites drive iteration and blocking defects; \textbf{shadow} deployments compare policy variants on live traffic shapes without raising user risk; \textbf{red teaming} probes adaptive attacker models; \textbf{regression} gates ensure monotonic improvement on standardized bundles when policies, models, or detectors change. Release criteria combine statistical thresholds, worst-case caps, mandatory evidence artifacts, and operator sign-off rules.

\researchagenda{Evaluation frontier}{How should organizations jointly optimize safety, utility, cost, latency, and auditability under non-deterministic decoding and policy interventions with measurable confidence intervals?}

\keytakeaway{The \EvalFramework{} converts opinion into evidence: versioned suites, explicit contracts, and correlated telemetry that the \ControlPlane{} can defend in audits and releases.}

% =============================
% Chapter 6
% =============================
\chapterstarter{Governance, Deployment, and Research Roadmap}{
This chapter operationalizes the blueprint: evidence-centric governance, phased delivery, roles, and a concise research agenda. \textbf{Pillar focus:} synthesis of \ControlPlane{}, \DecoupledSafety{}, and \EvalFramework{}. It closes the loop from engineering to assurance culture.
}{
\begin{itemize}
  \item Roadmap phases with deliverables and acceptance anchors.
  \item Operating model unifying science, runtime kernel, evidence/\textsc{ci}/release, and multilingual intelligence.
  \item Research themes aligned to compositional runtime safety and assurance under uncertainty.
\end{itemize}
}

\section{Evidence graph and release governance}
\textbf{Pillar:} \EvalFramework{} / \ControlPlane{}. Release decisions rely on an evidence graph linking benchmark runs, policy versions, model routes, detector versions, representative telemetry slices, and human approvals. This graph must support differential updates: partial reruns when only subsystems change. Governance integrates legal, privacy, and safety stakeholders through explicit acceptance contracts rather than ad hoc meetings alone.

\section{Organizational operating model}
\textbf{Pillar:} synthesis. Effective programs align four recurring capabilities: (i) security science and benchmark governance; (ii) runtime \SafetyKernel{} engineering; (iii) evidence, \textsc{ci}/\textsc{cd}, and release gating; (iv) multilingual and locale-specific intelligence as a source of semantic tests and detectors rather than as an afterthought. The objective is co-evolution of theory, systems, benchmarks, and operations---not isolated teams optimizing local metrics.

\section{Phased deployment roadmap}
\begin{itemize}
  \item \textbf{Phase 1 -- Baseline envelope.} Input/output \Guardrails{}, structured audit logging, minimal acceptance contracts and regression suites for highest-risk scenarios.
  \item \textbf{Phase 2 -- Contextual risks.} Retrieval mediation, tool-use policies, expanded evidence correlation, shadow evaluations.
  \item \textbf{Phase 3 -- Unified control plane.} Cross-model policy abstractions, centralized policy lifecycle, integrated dashboards for decisions and evidence.
  \item \textbf{Phase 4 -- Adaptive assurance.} Risk-adaptive selection with formal targets, stronger automation in regression discovery, targeted verification experiments where feasible.
\end{itemize}

\begin{deliverabletable}
Benchmark constitution & Documented suites + aggregation rules & CI gates pass with signed artifacts & Assurance lead & P1 \\
Acceptance contracts & Scenario thresholds + exceptions & Release board uses versioned bundles & Risk + Product & P1 \\
Safety kernel MVP & Event types + enforce + log & Measured reduction in bypass vs.\ baseline & Runtime & P1 \\
Retrieval + tool policies & Mediation rules live in prod & Incidents down vs.\ red-team catalog & Platform & P2 \\
Evidence service & Correlation IDs, immutability & Incident replay under time budget & SRE & P2 \\
\end{deliverabletable}

\section{Research roadmap}
\textbf{Pillar:} synthesis. Priority themes include compositional safety under heterogeneous stacks; runtime controllability with bounded regret trade-offs; policy--model separation with validated interfaces; evaluation under non-determinism via statistical and bayesian testing; security--utility--latency multi-objective optimization; and assurance science linking benchmarks to formal claims where possible.

\keytakeaway{A mature program ships controls, evidence, and learning together: engineering enforcement, measuring it honestly, and compounding organizational capability through reproducible incidents and benchmarks.}

% =============================
% Appendices
% =============================
\appendix

\chapter{Geo-specific and Multilingual Evaluation}
Regional obligations and languages introduce distinct risk semantics that should not be collapsed into a single global score. Practical patterns include: localized suite libraries with jurisdiction labels; parallel tracks sharing execution infrastructure but differing in categories, aggregation, and report templates; language-aware detectors and tokenizer pitfalls; mixed-language inputs; and explicit mapping from regulatory categories to test obligations. Evidence packs name the track, suite versions, and decision rationale for inclusion in unified release boards.

\chapter{Terminology and Role Boundaries}
\textbf{Security judgment provider} supplies policy interpretation, risk scoring, detection artifacts, explainability payloads, and structured evidence suitable for audits. \textbf{Execution platform} delivers networking, scaling, orchestration \textsc{sla}, and workload management. Boundaries prevent ambiguous ownership where neither side produces replayable traces. External names for internal programs map to these neutral roles in external communications.

\end{document}
